% =====================================================================
% Preambolo -- Manuale CSLR
% =====================================================================
\usepackage[T1]{fontenc}
\usepackage[utf8]{inputenc}
\usepackage[italian]{babel}
\usepackage{lmodern}
\usepackage[a4paper,margin=2.5cm]{geometry}
\usepackage{microtype}
\usepackage{xcolor}
\usepackage{enumitem}
\usepackage{booktabs}
\usepackage{array}
\usepackage{tabularx}
\usepackage{longtable}
\usepackage{amsmath}
\usepackage[most]{tcolorbox}
\usepackage{tikz}
\usetikzlibrary{positioning,arrows.meta}
\usepackage{hyperref}
\hypersetup{
  colorlinks=true,
  linkcolor=blue!50!black,
  urlcolor=blue!60!black,
  pdftitle={Cybersecurity Laws and Regulations -- Manuale del corso}
}

\setlist{itemsep=2pt,topsep=4pt}
\renewcommand{\arraystretch}{1.25}

% ---------------------------------------------------------------------
% Riquadri
% ---------------------------------------------------------------------
% Punto che il docente ha indicato come rilevante per l'esame
\newtcolorbox{esame}{breakable,colback=red!4,colframe=red!60!black,
  fonttitle=\bfseries,title=Per l'esame}

% Nota generica (titolo opzionale)
\newtcolorbox{nota}[1][Nota]{breakable,colback=gray!6,colframe=gray!60!black,
  fonttitle=\bfseries,title=#1}

% Informazione non confermata in aula o calcolata da noi: da verificare
\newtcolorbox{verifica}{breakable,colback=orange!6,colframe=orange!80!black,
  fonttitle=\bfseries,title=Da verificare}

% Soluzione di quiz / simulazioni (titolo opzionale)
\newtcolorbox{soluzione}[1][Soluzione]{breakable,colback=green!4,colframe=green!45!black,
  fonttitle=\bfseries,title=#1}

% Segnaposto per parti ancora da scrivere
\newcommand{\dacompletare}[1]{%
  \begin{tcolorbox}[colback=yellow!8,colframe=yellow!60!black,
    fonttitle=\bfseries,title=Da completare]#1\end{tcolorbox}}

% ---------------------------------------------------------------------
% Intestazione di una lezione
%   \lezione{data}{titolo}{docenti}{argomenti}
% ---------------------------------------------------------------------
\newcommand{\lezione}[4]{%
  \chapter[#1 -- #2]{#2}%
  \begin{tcolorbox}[colback=blue!3,colframe=blue!50!black,
    fonttitle=\bfseries,title=Scheda lezione]
  \textbf{Data:} #1\\
  \textbf{Docenti:} #3\\
  \textbf{Argomenti:} #4
  \end{tcolorbox}}

% ---------------------------------------------------------------------
% Quiz: domanda numerata + opzioni (a), (b), ...
% ---------------------------------------------------------------------
\newcounter{domanda}
\newenvironment{domanda}[1][]{%
  \refstepcounter{domanda}\par\medskip\noindent
  \textbf{Domanda \thedomanda}\ifx&#1&\else\ \textit{(#1)}\fi\textbf{.}~}%
  {\par}
\newenvironment{opzioni}{\begin{enumerate}[label=(\alph*),nosep,leftmargin=2.2em]}{\end{enumerate}}
